%% ch09-security-ref.tex — Security Reference
%% JA4proxy Technical Reference Manual

\chapter{Security Reference}
\index{security}
\label{ch:security}

\begin{chapterquote}
\ja\ is a security tool, but it is also a target. This chapter documents the security
properties of \ja\ itself: what it guarantees, where its boundaries are, what its
known gaps are at POC stage, and how to harden it for production.
\end{chapterquote}

\newthought{The security model of \ja} rests on three guarantees: it never decrypts
traffic (cryptographic transparency), it fails open when uncertain (false positive
minimisation), and it defaults to monitor mode (dial=0) so nothing is ever blocked
until the operator explicitly enforces. Browser protection comes from these defaults,
not from an unconditional ALPN bypass.

\section{The Bypass System — Security Model}
\index{bypass!security model}
\index{bypass!risk analysis}

\subsection{Bypass Architecture}

Bypass conditions are not simply configuration flags — they represent a deliberate
security policy decision about which classes of traffic should be exempt from scoring.
Each bypass type has a specific threat model and risk profile when enabled or disabled.

\begin{fullwidth}
\begin{table}[h]
\caption{Bypass security analysis — each bypass, its threat model, and risk of disabling}
\label{tab:bypass-security}
\begin{tabular}{@{}lp{3.5cm}p{3.5cm}p{4cm}@{}}
\toprule
\textbf{Bypass} & \textbf{Threat addressed} & \textbf{Risk if disabled} & \textbf{Use case for disabling} \\
\midrule
\code{alpn\_browser\_bypass} & Not applicable — this is a false-positive protection, not a threat bypass & Browser traffic scored; significant false positive risk at scale & Scoring specific h2 API clients known to not be real browsers \\
\addlinespace
\code{ja4\_whitelist\_bypass} & Trusted fingerprints (known-good tools, monitoring) bypass scorer & Whitelisted fingerprints may be blocked if they score above threshold & Testing whether a whitelisted fingerprint would be scored \\
\addlinespace
\code{mtls\_bypass} & Cryptographic identity proofing — mTLS client is verified & mTLS clients scored and potentially blocked & Auditing mTLS client behaviour \\
\addlinespace
\code{static\_ip\_allowlist} & Trusted internal IPs (monitoring, CI) bypass scorer & Monitoring may be blocked; CI tests fail & Temporary removal during security audit \\
\addlinespace
\code{ja4\_blacklist\_bypass} & Known-bad fingerprints (C2, scanners) immediately rejected & Blacklisted fingerprints only blocked if score exceeds dial threshold & Investigating traffic from a specific fingerprint \\
\addlinespace
\code{country\_blacklist\_bypass} & Geographically restricted access & Blocked countries go through scorer; may not be blocked at low dial & Investigating specific country traffic \\
\addlinespace
\code{spamhaus\_bypass} & Hijacked/criminal IP blocks immediately rejected & Spamhaus matches produce signal(+80) instead of hard block & Investigating Spamhaus-listed range \\
\addlinespace
\code{tls\_version\_bypass} & Deprecated TLS versions immediately rejected & Old TLS produces risk signal instead of hard reject & Auditing legacy client population \\
\bottomrule
\end{tabular}
\end{table}
\end{fullwidth}

\subsection{Startup Warnings for High-Risk Bypass Disabling}
\index{bypass!startup warnings}

When any bypass is disabled that has elevated false positive or false negative risk,
\ja\ emits a \code{WARN} level log at startup. This ensures the operator cannot
accidentally disable a bypass in a config file and have the change go unnoticed.

\begin{lstlisting}[language=bash]
# Example startup warnings for disabled bypasses
WARN | policy | event=bypass_disabled | bypass=alpn_browser_bypass
     | effect=browser traffic will be scored; false positive risk elevated

WARN | policy | event=bypass_disabled | bypass=spamhaus_drop_bypass
     | effect=Spamhaus DROP matches scored as signal(+80) instead of hard block
\end{lstlisting}

\section{The Asymmetry Principle — In Depth}
\index{asymmetry principle!in depth}
\index{false positive!protection}

\subsection{Why Fail Open is Correct}

The asymmetry principle (Chapter~\ref{ch:architecture}) has specific technical
manifestations throughout the system. Understanding why each mechanism exists as it
does requires understanding the cost model:

\begin{itemize}
  \item \textbf{A false negative} (bad bot allowed) produces a malformed request to the
    backend. The backend logs it. The analytics node may detect the pattern retrospectively.
    The next connection from the same IP will have a better signal set (AbuseIPDB cache
    populated, beaconing window extended). Recovery is automatic.
  \item \textbf{A false positive} (real browser blocked) produces a failed page load for
    a human user. The user sees a connection error with no explanation. They may never
    return. There is no automatic recovery — the damage is immediate and irreversible.
\end{itemize}

\subsection{False Positive Protection Mechanisms}

These mechanisms exist specifically to prevent false positives:

\begin{description}
  \item[ALPN browser bypass] An optional (disabled-by-default) bypass that allows h2/h1
    ALPN connections without scoring. It is \emph{disabled by default} because bots can
    trivially set ALPN values to impersonate browsers. When disabled, browser-looking
    traffic is scored like any other connection. Real browser protection comes from
    monitor mode (dial=0), fail-open defaults, and whitelisting known-good JA4
    fingerprints, not from an unconditional ALPN bypass.\sidenote{If the operator
    enables the bypass, they accept the risk that bots can use ALPN spoofing to bypass
    scoring.}

  \item[ALLOW cached 30min, BLOCK cached 30s] The local LRU cache intentionally uses
    asymmetric TTLs. An IP allowed 29 minutes ago gets another 30-minute free pass.
    An IP blocked 31 seconds ago is re-evaluated against current Redis state.

  \item[Local cache wins on Redis failure] When the local cache says allow and Redis
    says block (or Redis is unreachable), the local cache wins. Real users in the local
    cache continue to be served even during Redis outages.

  \item[External service unavailability → zero signal] If AbuseIPDB is down, RDAP is
    unreachable, or DNS lookups time out, the affected signal contributes 0 to the score.
    The score can only go down when services fail, not up.

  \item[AbuseIPDB floor at 50] AbuseIPDB scores below 50 are suppressed. This prevents
    shared infrastructure (NAT gateways, corporate egress, VPNs) from triggering signals
    just because one user behind that NAT misbehaved.

  \item[RDAP block expansion off by default] Expanding a block to an entire /24 based
    on one bad IP is a high-impact action. The default is off. Operators must consciously
    opt in.
\end{description}

\section{The Fail-Open Guarantee}
\index{fail open!guarantee}

The fail-open guarantee is: \textbf{in all uncertainty, connections are allowed rather
than blocked}. This guarantee applies to:

\begin{itemize}
  \item Redis unreachable: bans not enforced; rate limits not enforced; signals that
    depend on Redis return 0
  \item External service (AbuseIPDB, RDAP, DNS) unreachable: affected signals return 0
  \item Config parse error on reload: previous valid config retained; new config not applied
  \item Pipeline exception: connection is allowed; exception is logged
  \item Unknown fingerprint: score 0 (unknown is not inherently suspicious)
\end{itemize}

\begin{importantbox}[Fail-Open is Not a Bug]
Operators sometimes ask whether fail-open can be changed to fail-closed. The answer is
no — this would be an architectural change that fundamentally compromises the false
positive guarantee. Fail-closed means an outage blocks real users. The system is
designed to tolerate outages gracefully, not to use them as an excuse to block.
\end{importantbox}

\section{Hot Config Reload — Security Implications}
\index{hot reload!security implications}

Hot config reload is a convenience, but it has security implications:

\begin{description}
  \item[Policy change audit log] Every config reload is written to
    \code{management:policy\_audit} in Redis with a timestamp and attribution.
    Operators can audit all policy changes without consulting application logs.
  \item[Config file integrity] The proxy does not verify the cryptographic integrity
    of \code{config/proxy.yml} before loading. In production, config files should
    be managed by a configuration management system (Ansible, Terraform, Kubernetes
    ConfigMap) with change control and audit trail.
  \item[SIGHUP race condition] There is a brief window during config reload where
    new bypass states are being applied and some connections may be evaluated under
    the old configuration. This window is bounded by the duration of the reload
    operation (typically $<$50ms). No connections are dropped during reload.
  \item[Management UI reload] The Management UI triggers reload via a Redis pub/sub
    message. If Redis is compromised, an attacker could trigger arbitrary config
    reloads. Redis authentication (\configkey{redis.password}) must be set in
    production.
\end{description}

\section{Phase 200-Series Hardening}
\index{Phase 200-series}
\index{hardening!Phase 200-series}
\label{sec:phase200-hardening}

\newthought{The Phase 200-series closes the production-hardening gaps} flagged
in Phase~14 and the comprehensive security audit. These items are now in
production builds of the Go proxy.

\subsection{Redis TLS (Phase 201)}
\index{Redis!TLS}

Redis communications are now TLS-encrypted by default. The Go proxy connects
to Redis with mutual TLS where the Redis instance is configured with
\code{tls-port} and a server certificate; client certificates are loaded from
the path configured under \code{redis.tls.client\_cert\_path}. Plaintext Redis
connections remain available for development (a Unix-socket-only configuration
with restrictive filesystem permissions is the recommended local-dev path), but
production Helm charts default to TLS-on. See \code{docs/phases/PHASE\_201.md}
for the full configuration matrix.

\subsection{PROXY Protocol v2 with TLVs (Phase 203)}
\index{PROXY protocol!v2}
\index{PROXY protocol!TLVs}

The Go proxy now parses PROXY protocol v2 binary headers, including TLV
(type-length-value) extensions. The TLV trust path is gated by
\code{security\_policy.trusted\_upstream\_cidrs}: a TLV is honoured only when
the TCP peer is within the configured trusted CIDR. This prevents a
direct-to-port-8080 attacker from forging TLVs to claim a whitelisted source
IP. The trust check runs before any TLV is parsed, so a malformed TLV from an
untrusted peer cannot trigger a parser path at all.

\subsection{Default Credentials Removed (Phase 202)}
\index{credentials!default removed}

Default credentials have been removed from the Docker image, Helm chart, and
\code{docker-compose.yml}. Previously, Redis shipped with no \code{requirepass},
Grafana shipped with \code{admin/admin}, and the Management API had a hard-coded
local-dev token. All three now require explicit env-var provisioning at startup;
the proxy refuses to start if a required credential is unset. The Helm chart
fails its preflight check and Docker Compose surfaces a clear error rather than
silently starting with a blank password. See \code{docs/phases/PHASE\_202.md}
for the env-var reference and the migration runbook.

\section{Known Security Gaps (POC Stage)}
\index{security gaps!POC stage}

The following gaps are documented in \code{docs/DMZ\_DEPLOYMENT\_READINESS.md} and
\code{docs/security/COMPREHENSIVE\_SECURITY\_AUDIT.md}. They must be addressed
before production deployment with dial $>$ 0.

\begin{fullwidth}
\begin{table}[h]
\caption{Known security gaps at POC stage}
\label{tab:security-gaps}
\begin{tabular}{@{}lllp{5cm}@{}}
\toprule
\textbf{Gap} & \textbf{Severity} & \textbf{Phase addressing it} & \textbf{Description} \\
\midrule
Secrets in environment variables & Medium & Phase 14 & API keys and Redis password should come from a secret manager (Vault, AWS Secrets Manager), not plain environment variables \\
Redis without TLS & High & Phase 14 & Redis communications are unencrypted. In production, Redis must use TLS or a Unix domain socket with filesystem permissions \\
Tarpit self-protection absent & Medium & Phase 14 & The tarpit mechanism has no limit on concurrent tarpitted connections. A sustained flood could exhaust file descriptors \\
Config file not integrity-checked & Low & Phase 14 & No signature verification on proxy.yml before hot-reload \\
Management UI authentication & High & Phase 13 & Management UI has no authentication in POC. Must add before production \\
Redis ACLs not configured & Medium & Phase 34 & All clients share a single Redis connection. Should use ACLs with per-service credentials \\
No Seccomp/AppArmor profiles & Low & Phase 55 & Container syscall surface not restricted \\
Supply chain integrity & Low & Phase 35 & No SBOM, no container image signing \\
\bottomrule
\end{tabular}
\end{table}
\end{fullwidth}

\section{Production Hardening Checklist}
\index{production hardening!checklist}
\index{deployment!security hardening}

\begin{enumerate}
  \item[\(\square\)] \textbf{Secrets manager integration} — API keys and passwords from Vault or cloud secret manager; never in \code{.env} files
  \item[\(\square\)] \textbf{Redis TLS} — enable TLS for Redis connections or use Unix socket with restricted permissions
  \item[\(\square\)] \textbf{Redis AUTH} — set strong Redis password; verify it is not default
  \item[\(\square\)] \textbf{Redis ACLs} — separate credentials for proxy, analytics, and management services
  \item[\(\square\)] \textbf{Management UI authentication} — complete Phase 13 before deploying Management UI
  \item[\(\square\)] \textbf{Tarpit resource limits} — configure maximum concurrent tarpitted connections
  \item[\(\square\)] \textbf{Container non-root user} — verify proxy runs as non-root (uid 1000)
  \item[\(\square\)] \textbf{Read-only filesystem} — mount application directories read-only where possible
  \item[\(\square\)] \textbf{Network policy} — enforce network zone isolation (see Chapter~\ref{ch:deployment})
  \item[\(\square\)] \textbf{HAProxy trusted CIDR} — configure the exact HAProxy IP range in proxy trusted upstream config
  \item[\(\square\)] \textbf{Grafana admin password changed} — default \code{admin}/\code{admin} must not be used
  \item[\(\square\)] \textbf{Loki log retention} — configure retention policy (30 days recommended for GDPR compliance)
  \item[\(\square\)] \textbf{Audit log retention} — verify \code{management:policy\_audit} retention policy
  \item[\(\square\)] \textbf{Pentest validation} — complete the adversarial test suite (Phase 45) before raising dial
\end{enumerate}

\section{IP Spoofing Protection}
\index{IP spoofing}
\index{PROXY protocol!trust}

Without trusted upstream CIDR verification, an attacker who can reach port 8080
directly could inject a PROXY protocol header claiming any source IP address —
effectively bypassing all IP-based checks by claiming to be a whitelisted IP.

\ja\ mitigates this by verifying that the sender of the TCP connection (before reading
the PROXY protocol header) is within the configured trusted upstream CIDR. Only
HAProxy addresses should be in this CIDR. If the sender is not in the trusted CIDR,
the PROXY protocol header is rejected and the real TCP peer IP is used instead.

\begin{warningbox}[Port 8080 Must Not Be Internet-Accessible]
Port 8080 (the \ja\ listen port) must only be reachable from HAProxy. If an attacker
can reach port 8080 directly, they can manipulate their apparent source IP. Firewall
rules must block port 8080 from all sources except the HAProxy IP range.
\end{warningbox}

\section{JA4 Fingerprint Limitations}
\index{JA4!limitations}
\index{fingerprint evasion}

JA4 fingerprinting is powerful but not infallible. Security engineers should
understand its limitations:

\begin{description}
  \item[JA4 copying] An attacker can copy a known-good JA4 fingerprint by configuring
    their TLS library to match the exact cipher suite list and extension ordering.
    JA4T (TCP fingerprint) makes this harder — it requires matching the OS-level TCP
    stack, not just the TLS library.
  \item[JA4 rotation] A sophisticated attacker can rotate through JA4 values to evade
    per-fingerprint rate limits. The \code{by\_ip\_ja4\_pair} rate limit strategy is
    designed to catch this pattern.
  \item[JA4 collisions] Multiple legitimate clients may share the same JA4 value.
    JA4 blacklisting should only be applied to fingerprints that are \textit{uniquely}
    associated with malicious tools — not fingerprints shared with any legitimate client
    software.
\end{description}
